% Generated draft: requires mathematical and source review.
\documentclass[11pt]{article}
\usepackage[margin=25mm]{geometry}
\usepackage{fontspec}
\setmainfont{Latin Modern Roman}
\usepackage{amsmath,amssymb,mathtools,unicode-math}
\setmathfont{Latin Modern Math}
\usepackage{enumitem,needspace,xurl,hyperref}
\setlength{\parindent}{0pt}
\setlength{\parskip}{5pt}
\setlength{\emergencystretch}{3em}
\newcommand{\N}{\mathbb{N}}
\newcommand{\Z}{\mathbb{Z}}
\newcommand{\Q}{\mathbb{Q}}
\newcommand{\R}{\mathbb{R}}
\newcommand{\C}{\mathbb{C}}
\newcommand{\F}{\mathbb{F}}
\newcommand{\A}{\mathbb{A}}
\newcommand{\PP}{\mathbb{P}}
\newcommand{\E}{\mathbb{E}}
\newcommand{\eps}{\varepsilon}
\newcommand{\dd}{\,\mathrm{d}}
\begin{document}
\section*{UnsolvedMath ID: 30006433}
\textbf{Source problem code:} catalogue entry
\subsection{Definitions and mathematical statement}
Let \(A\) range over the algebras intended by the supplied source's phrase ``the subalgebra-zeta setting,'' and let \(d\) denote their dimension in the sense intended there. Write \(Z_A^{\mathrm{top}}(s)\) for the topological subalgebra zeta function attached to \(A\). The supplied record does not specify the base field or ring, the admissible class of algebras, the meaning of dimension over that base, or the construction and normalization of this zeta function. Thus this notation records the intended invariant without supplying its missing definition.

For a nonzero rational function \(R(s)=P(s)/Q(s)\), where \(P\) and \(Q\) are nonzero polynomials over its coefficient field, define
\[
\deg R=\deg P-\deg Q.
\]
Equivalently, its order of vanishing at infinity is \(\operatorname{ord}_{t=0}R(1/t)=-\deg R\). This degree convention requires the intended zeta function to be a nonzero rational function; its construction and these prerequisites must be checked in the underlying source.

\textbf{Degree conjecture (source reference [8, Conjecture I]).} For every algebra \(A\) in the intended admissible class, with dimension parameter \(d\),
\[
\deg Z_A^{\mathrm{top}}(s)=-d.
\]
Under the rational-function interpretation above, this says that \(Z_A^{\mathrm{top}}(s)\) has order of vanishing \(d\) at infinity. The missing hypotheses defining the admissible class prevent this transcription from being a fully specified universal mathematical statement.

\textbf{Qualifications from the supplied context.} The record also discusses a reduced subalgebra zeta function \(Z_A^{\mathrm{red}}(T)\), whose definition is likewise absent. It reports that Evseev [6, Proposition 4.1], under a particular basis hypothesis that is not reproduced, constructs a \(d\)-dimensional cone \(C\subseteq\mathbb R_{\geq0}^{d}\) for which \(Z_A^{\mathrm{red}}(T)\) is the coarse Hilbert series of \(\mathbb Q[C\cap\mathbb Z^d]\). Here the monoid algebra has basis \(e^m\) indexed by \(m\in C\cap\mathbb Z^d\), with multiplication \(e^m e^n=e^{m+n}\); the grading used for the stated Hilbert series is not supplied. The excerpt further reports that Evseev's explicit cone description, together with results cited as [11, Chapter 1], makes this algebra Gorenstein. These are supplied contextual claims under an omitted basis hypothesis, not assertions established by this draft.

The adjacent Hilbert series conjecture, attributed to Voll, asks for a ``natural, meaningful'' \(\mathbb N_0\)-graded Gorenstein algebra of dimension \(d\) whose Hilbert series is \(Z_A^{\mathrm{red}}(T)\), where \(\mathbb N_0=\{0,1,2,\ldots\}\). For a graded algebra \(B=\bigoplus_{n\geq0}B_n\) over a field \(k\), with finite-dimensional graded pieces, its Hilbert series means \(\sum_{n\geq0}(\dim_k B_n)T^n\). The field, dimension convention, grading conditions, and meaning of ``natural, meaningful'' in the conjecture are not specified in the record. The source qualifies this conjecture as informally circulated for some time, with more cautious published formulations. It also mentions a coincidence conjecture relating reduced and topological zeta functions, but supplies no formula for it. These neighboring conjectures are context; the requested target is the degree conjecture above.
\subsection{Short English statement}
The supplied conjecture predicts that the topological subalgebra zeta function of every admissible algebra of dimension \(d\) has rational-function degree \(-d\), equivalently vanishing order \(d\) at infinity. The admissible algebras, base, dimension convention, and zeta-function construction are missing, so the statement remains provisional. The supplied status is open; neither that status nor the contextual results have been independently verified.
\subsection{Sources}
\begin{itemize}
\item Supplied local UnsolvedMath registry record: permanent ID \(30006433\), title ``Degree of Topological Subalgebra Zeta Functions,'' exact source problem code \texttt{catalogue entry}, in the Oberwolfach Reports Open Problems collection. No proposer is named. The supplied proposed year is 2025.
\item Supplied report citation: \emph{Cohomology Theories for Automorphic Forms and Enumerative Algebra} (2026), Oberwolfach Reports, pp. 2065--2124, DOI \url{https://doi.org/10.4171/owr/2025/38}. The citation year 2026 and the supplied 2025 report/proposal metadata require bibliographic reconciliation.
\item Supplied literature source: ``Oberwolfach Report 38,'' \url{https://ems.press/content/serial-article-files/52251}.
\item Supplied literature source: ``Rossmann--Voll survey slides,'' \url{https://torossmann.github.io/files/mfo25b.pdf}.
\item The literature assessment dated \texttt{2026-08-22} is supplied dated curation, not independently verified here. It labels the problem open and states that the degree assertion is still presented as a conjecture, while known cone and Hilbert-series constructions cover special bases and classes rather than all algebras in the extracted scope. The supplied background explicitly requests independent verification of time-sensitive status and cited claims.
\item Internal references retained from the excerpt: [8, Conjecture I] for the degree conjecture; Evseev [6, Proposition 4.1] for the special-basis cone construction; [11, Chapter 1] for the cited commutative-algebra results; and Voll for the adjacent Hilbert series conjecture. The bibliography entries resolving [6], [8], and [11] are not supplied, so their authorship and publication details cannot be completed from this record.
\end{itemize}
\end{document}
